\documentclass[12pt, a4paper]{book}

% ─── Core packages ────────────────────────────────────────────────────────────
\usepackage[T1]{fontenc}
\usepackage[utf8]{inputenc}
\usepackage[charter, cal=cmcal]{mathdesign}
\usepackage[scaled=0.92]{helvet}
\usepackage{geometry}
\geometry{
  a4paper,
  top=30mm, bottom=30mm,
  inner=32mm, outer=25mm,
  marginparwidth=18mm
}

% ─── Mathematics ──────────────────────────────────────────────────────────────
\usepackage{amsmath, amsthm}

% ─── Chemistry ────────────────────────────────────────────────────────────────
\usepackage[version=4]{mhchem}

% ─── Graphics & colour ────────────────────────────────────────────────────────
\usepackage{graphicx}
\usepackage[dvipsnames,svgnames,table]{xcolor}
\usepackage{tikz}
\usetikzlibrary{arrows.meta, shapes.geometric, positioning,
                calc, backgrounds, decorations.pathmorphing,
                decorations.markings, patterns}
\usepackage{pgfplots}
\pgfplotsset{compat=1.18}
\usepackage{tcolorbox}
\tcbuselibrary{skins, breakable, theorems}

% ─── Layout ───────────────────────────────────────────────────────────────────
\usepackage{titlesec}
\usepackage{fancyhdr}
\usepackage{enumitem}
\usepackage{booktabs}
\usepackage{array}
\usepackage{multicol}
\usepackage{marginnote}
\usepackage{mdframed}
\usepackage{wrapfig}
\usepackage{caption}
\usepackage{subcaption}
\usepackage{parskip}
\usepackage{siunitx}
\usepackage[colorlinks=true, linkcolor=BrickRed,
            urlcolor=NavyBlue, citecolor=ForestGreen]{hyperref}

% ══════════════════════════════════════════════════════════════════════════════
% COLOUR PALETTE  (identical to Module 1.1)
% ══════════════════════════════════════════════════════════════════════════════
\definecolor{JourneyColor}{HTML}{065F46}   % forest green — journey headings
\definecolor{ModuleColor}{HTML}{0F766E}    % deep teal — module headings
\definecolor{PuzzleColor}{HTML}{B45309}    % deep amber — puzzle boxes
\definecolor{DiscoveryColor}{HTML}{D97706} % golden amber — discovery boxes
\definecolor{WarningColor}{HTML}{E11D48}   % coral red — "model breaks" boxes
\definecolor{ProblemColor}{HTML}{0F766E}   % dark teal — problem sets
\definecolor{ReflectColor}{HTML}{047857}   % green — reflection boxes
\definecolor{HistoryColor}{HTML}{9A3412}   % brown — history notes
\definecolor{LightBg}{HTML}{F8FAFC}        % near-white background

% ══════════════════════════════════════════════════════════════════════════════
% TCOLORBOX ENVIRONMENTS
% ══════════════════════════════════════════════════════════════════════════════
\newtcolorbox{puzzlebox}[1][]{
  enhanced, breakable,
  colback=PuzzleColor!6, colframe=PuzzleColor,
  fonttitle=\bfseries\sffamily, title={The Puzzle},
  attach boxed title to top left={yshift=-2mm, xshift=6mm},
  boxed title style={colback=PuzzleColor, colframe=PuzzleColor},
  borderline west={3pt}{0pt}{PuzzleColor}, #1}

\newtcolorbox{discoverybox}[1][]{
  enhanced, breakable,
  colback=DiscoveryColor!8, colframe=DiscoveryColor,
  fonttitle=\bfseries\sffamily, title={Discovery},
  attach boxed title to top left={yshift=-2mm, xshift=6mm},
  boxed title style={colback=DiscoveryColor, colframe=DiscoveryColor},
  borderline west={3pt}{0pt}{DiscoveryColor}, #1}

\newtcolorbox{problembox}[1][]{
  enhanced, breakable,
  colback=ProblemColor!5, colframe=ProblemColor,
  fonttitle=\bfseries\sffamily, title={Problems},
  attach boxed title to top left={yshift=-2mm, xshift=6mm},
  boxed title style={colback=ProblemColor, colframe=ProblemColor},
  borderline west={3pt}{0pt}{ProblemColor}, #1}

\newtcolorbox{deepbox}[1][]{
  enhanced, breakable,
  colback=JourneyColor!5, colframe=JourneyColor,
  fonttitle=\bfseries\sffamily, title={Deep Problems},
  attach boxed title to top left={yshift=-2mm, xshift=6mm},
  boxed title style={colback=JourneyColor, colframe=JourneyColor},
  borderline west={4pt}{0pt}{JourneyColor}, #1}

\newtcolorbox{historybox}[1][]{
  enhanced, breakable,
  colback=HistoryColor!7, colframe=HistoryColor!60,
  fonttitle=\bfseries\sffamily\small, title={Historical Note},
  attach boxed title to top left={yshift=-2mm, xshift=6mm},
  boxed title style={colback=HistoryColor!80, colframe=HistoryColor}, #1}

\newtcolorbox{reflectionbox}[1][]{
  enhanced, breakable,
  colback=ReflectColor!6, colframe=ReflectColor,
  fonttitle=\bfseries\sffamily, title={Reflection},
  attach boxed title to top left={yshift=-2mm, xshift=6mm},
  boxed title style={colback=ReflectColor, colframe=ReflectColor},
  borderline west={3pt}{0pt}{ReflectColor}, #1}

\newtcolorbox{modelbreaks}[1][]{
  enhanced, breakable,
  colback=WarningColor!6, colframe=WarningColor,
  fonttitle=\bfseries\sffamily, title={Where This Model Breaks},
  attach boxed title to top left={yshift=-2mm, xshift=6mm},
  boxed title style={colback=WarningColor, colframe=WarningColor},
  borderline west={4pt}{0pt}{WarningColor}, #1}

\newtcolorbox{namedconcept}[2][]{
  enhanced,
  colback=ModuleColor!6, colframe=ModuleColor,
  fonttitle=\bfseries\sffamily\large, title={#2},
  attach boxed title to top left={yshift=-2mm, xshift=6mm},
  boxed title style={colback=ModuleColor, colframe=ModuleColor}, #1}

% ══════════════════════════════════════════════════════════════════════════════
% THEOREM ENVIRONMENTS
% ══════════════════════════════════════════════════════════════════════════════
\theoremstyle{definition}
\newtheorem{observation}{Observation}[chapter]
\newtheorem{question}{Question}[chapter]
\newtheorem{experiment}{Experiment}[chapter]

% ══════════════════════════════════════════════════════════════════════════════
% SECTION FORMATTING
% ══════════════════════════════════════════════════════════════════════════════
\titleformat{\chapter}[display]
  {\normalfont\sffamily\huge\bfseries\color{JourneyColor}}
  {\textcolor{JourneyColor!60}{\chaptertitlename\ \thechapter}}
  {10pt}{\Huge}[\vspace{2pt}\textcolor{JourneyColor}{\titlerule[1.5pt]}]

\titleformat{\section}
  {\normalfont\sffamily\Large\bfseries\color{ModuleColor}}
  {\thesection}{1em}{}[\textcolor{ModuleColor!40}{\titlerule[0.6pt]}]

\titleformat{\subsection}
  {\normalfont\sffamily\large\bfseries\color{JourneyColor!80}}
  {\thesubsection}{1em}{}

\titleformat{\subsubsection}
  {\normalfont\sffamily\normalsize\bfseries}
  {\thesubsubsection}{1em}{}

% ══════════════════════════════════════════════════════════════════════════════
% HEADER / FOOTER
% ══════════════════════════════════════════════════════════════════════════════
\pagestyle{fancy}
\fancyhf{}
\fancyhead[LE]{\small\sffamily\textcolor{JourneyColor}{\leftmark}}
\fancyhead[RO]{\small\sffamily\textcolor{ModuleColor}{\rightmark}}
\fancyfoot[C]{\small\thepage}
\renewcommand{\headrulewidth}{0.4pt}

% ══════════════════════════════════════════════════════════════════════════════
% CUSTOM COMMANDS
% ══════════════════════════════════════════════════════════════════════════════
\newcommand{\chem}[1]{\ce{#1}}
\newcommand{\important}[1]{\textcolor{ModuleColor}{\textbf{#1}}}
\newcommand{\newterm}[1]{\textit{\textcolor{DiscoveryColor}{#1}}}
\newcommand{\sidenote}[1]{%
  \marginnote{\footnotesize\sffamily\textcolor{JourneyColor!70}{#1}}}

% ══════════════════════════════════════════════════════════════════════════════
% DOCUMENT
% ══════════════════════════════════════════════════════════════════════════════
\begin{document}
\begin{namedconcept}{Law of Conservation of Mass}
In any chemical reaction,
the total mass of the reactants equals
the total mass of the products.

\begin{center}
\Large
$m_{\text{reactants}} = m_{\text{products}}$
\end{center}

Mass is neither created nor destroyed in chemical reactions.
It is \emph{rearranged}.
\end{namedconcept}

\vspace{0.5em}

This law sounds almost trivial today.
In 1785 it was a revolution.
It transformed chemistry from a qualitative art
into a quantitative science.

For the first time, chemists could \emph{predict}:
if you know the masses of all ingredients,
you know the total mass of all products.
You cannot get something from nothing.
You cannot make mass disappear.

\subsubsection{Applying the Law: Worked Examples}

\textbf{Example 1.} Iron reacts with sulfur to form iron(II) sulfide:
\begin{center}
\ce{Fe + S -> FeS}
\end{center}
A student mixes \SI{5.60}{\gram} of iron with excess sulfur
and heats strongly.
After reaction, the product weighs \SI{8.80}{\gram}.

How much sulfur was consumed?

\vspace{0.3em}
\textit{Solution.}
By conservation of mass, mass of sulfur consumed
$= \SI{8.80}{\gram} - \SI{5.60}{\gram} = \SI{3.20}{\gram}$.

\vspace{0.5em}

\textbf{Example 2.} Magnesium burns in air:
\begin{center}
\ce{2Mg + O2 -> 2MgO}
\end{center}
\SI{12.0}{\gram} of magnesium produces \SI{20.0}{\gram} of magnesium oxide.

How much oxygen was consumed?
What fraction of the product by mass is oxygen?

\vspace{0.3em}
\textit{Solution.}
Oxygen consumed $= \SI{20.0}{\gram} - \SI{12.0}{\gram} = \SI{8.0}{\gram}$.\\
Fraction that is oxygen $= 8.0/20.0 = 0.40$ (i.e.\ \SI{40}{\percent} by mass).

\begin{problembox}
\begin{enumerate}[label=\textbf{\arabic*.}]
  \item Calcium carbonate decomposes when heated:
        \begin{center}
        \ce{CaCO3 ->[\Delta] CaO + CO2}
        \end{center}
        \SI{50.0}{\gram} of \ce{CaCO3} produces \SI{28.0}{\gram} of \ce{CaO}.
        \begin{enumerate}[label=(\alph*)]
          \item What mass of \ce{CO2} was produced?
          \item Where did this \ce{CO2} go if the reaction happened in an open crucible?
          \item If a student only weighed the solid before and after
                (not capturing the gas), would they conclude mass was not conserved?
                What does this tell us about Lavoisier's methodology?
        \end{enumerate}

  \item Hydrogen gas burns in oxygen:
        \begin{center}
        \ce{2H2 + O2 -> 2H2O}
        \end{center}
        \SI{4.0}{\gram} of hydrogen completely reacts with oxygen.
        \begin{enumerate}[label=(\alph*)]
          \item Use conservation of mass to find the mass
                of water produced, given that \SI{32.0}{\gram} of \ce{O2} reacted.
          \item What is the mass ratio of hydrogen to oxygen in water?
          \item Would you expect this ratio to change if you used
                water from a different source?
        \end{enumerate}

  \item A sealed container holds \SI{10.0}{\gram} of wood and air.
        The wood burns completely inside the sealed container.
        The container is not opened.
        \begin{enumerate}[label=(\alph*)]
          \item What is the total mass of everything in the container
                after burning?
          \item If you then open the container, could the total mass change?
                Under what conditions?
        \end{enumerate}
\end{enumerate}
\end{problembox}

% ─── Section: What Reactions Are ─────────────────────────────────────────────
\subsection{What a Reaction Actually Is: Rearrangement}

The law of conservation of mass tells us something profound about
the nature of chemical reactions.

If mass is conserved, then reactions do not create or destroy matter.
They \emph{rearrange} it.
The atoms present before the reaction are the same atoms
present after — just connected differently.

This is a deep conceptual shift.

\begin{center}
\begin{tikzpicture}[
  every node/.style={font=\sffamily\small},
  atom/.style={circle, minimum size=0.7cm, font=\sffamily\footnotesize\bfseries}
]
  % Before reaction: H2 + Cl2
  \node[font=\sffamily\small\bfseries, color=JourneyColor]
        at (-1, 1.8) {Before:};

  % H2 molecule
  \node[atom, draw=blue!70, fill=blue!20] (h1) at (1.0, 1.8) {H};
  \node[atom, draw=blue!70, fill=blue!20] (h2) at (2.0, 1.8) {H};
  \draw[thick, blue!70] (h1) -- (h2);

  \node at (3.2, 1.8) {\large $+$};

  % Cl2 molecule
  \node[atom, draw=green!60!black, fill=green!20] (cl1) at (4.4, 1.8) {Cl};
  \node[atom, draw=green!60!black, fill=green!20] (cl2) at (5.4, 1.8) {Cl};
  \draw[thick, green!60!black] (cl1) -- (cl2);

  % Arrow
  \draw[-{Stealth[length=6pt]}, ultra thick, color=DiscoveryColor]
        (3.2, 0.9) -- (3.2, 0.1) node[midway, right]{react};

  % After reaction: 2 HCl
  \node[font=\sffamily\small\bfseries, color=JourneyColor]
        at (-1, -0.5) {After:};

  \node[atom, draw=blue!70, fill=blue!20] (h1b) at (1.0, -0.5) {H};
  \node[atom, draw=green!60!black, fill=green!20] (cl1b) at (2.0, -0.5) {Cl};
  \draw[thick, purple] (h1b) -- (cl1b);

  \node at (3.2, -0.5) {\large $+$};

  \node[atom, draw=blue!70, fill=blue!20] (h2b) at (4.4, -0.5) {H};
  \node[atom, draw=green!60!black, fill=green!20] (cl2b) at (5.4, -0.5) {Cl};
  \draw[thick, purple] (h2b) -- (cl2b);

  % Annotation
  \node[font=\sffamily\footnotesize\itshape, text=HistoryColor]
        at (3.2, -1.6) {Same atoms. New connections.};
\end{tikzpicture}
\end{center}

\vspace{0.5em}

This picture — atoms rearranging during reactions —
is not proven yet.
We have only established that mass is conserved.
But the rearrangement picture is the \emph{simplest hypothesis}
consistent with that fact.

We will confirm and deepen this picture dramatically
in Module~1.4, when we examine the ratios by mass
in which elements always combine.

% ─── Section: Types of Chemical Reaction ─────────────────────────────────────
\subsection{A First Classification of Reactions}

Now that we understand reactions as rearrangements,
we can start to classify them by the type of rearrangement.

\begin{center}
\begin{tikzpicture}[
  every node/.style={font=\sffamily\small, align=center},
  box/.style={draw=JourneyColor!60, fill=JourneyColor!8,
              rounded corners=3pt, minimum width=3.4cm,
              minimum height=1.2cm},
  subbox/.style={draw=ModuleColor!50, fill=ModuleColor!6,
                 rounded corners=2pt, minimum width=3.0cm,
                 minimum height=0.9cm},
  arr/.style={-{Stealth[length=4pt]}, thick, color=JourneyColor!70}
]
  \node[box, minimum width=4.5cm, minimum height=1cm]
        (react) at (0,0) {\textbf{Types of Chemical Reaction}};

  \node[subbox] (comb) at (-5.5,-2.2)
        {Combination\\$A + B \to AB$};
  \node[subbox] (decomp) at (-2,-2.2)
        {Decomposition\\$AB \to A + B$};
  \node[subbox] (displ) at (2,-2.2)
        {Single displacement\\$A + BC \to AC + B$};
  \node[subbox] (dd) at (5.5,-2.2)
        {Double displacement\\$AB + CD \to AD + CB$};

  \foreach \mynode in {comb, decomp, displ, dd}
    \draw[arr] (react) -- (\mynode);

  % Examples below each
  \node[font=\sffamily\footnotesize\itshape, text=HistoryColor]
        at (-5.5,-3.4) {\ce{2H2 + O2 -> 2H2O}};
  \node[font=\sffamily\footnotesize\itshape, text=HistoryColor]
        at (-2,-3.4) {\ce{2HgO -> 2Hg + O2}};
  \node[font=\sffamily\footnotesize\itshape, text=HistoryColor]
        at (2,-3.4) {\ce{Zn + 2HCl -> ZnCl2 + H2}};
  \node[font=\sffamily\footnotesize\itshape, text=HistoryColor]
        at (5.5,-3.4) {\ce{NaCl + AgNO3 -> AgCl + NaNO3}};
\end{tikzpicture}
\end{center}

\vspace{0.5em}

Each type has a characteristic pattern.
In a combination reaction, two simpler things join.
In decomposition, one thing breaks apart.
In displacement, one element displaces another from a compound.

\begin{problembox}
\begin{enumerate}[label=\textbf{\arabic*.}]
  \item Classify each reaction as combination, decomposition,
        single displacement, or double displacement:
        \begin{enumerate}[label=(\alph*)]
          \item \ce{2Na + Cl2 -> 2NaCl}
          \item \ce{CaCO3 -> CaO + CO2}
          \item \ce{Fe + CuSO4 -> FeSO4 + Cu}
          \item \ce{H2SO4 + 2KOH -> K2SO4 + 2H2O}
          \item \ce{N2 + 3H2 -> 2NH3}
        \end{enumerate}

  \item The reaction \ce{Fe + CuSO4 -> FeSO4 + Cu}
        shows iron displacing copper from solution.
        If you placed a strip of copper metal into iron sulfate solution,
        would you expect the reverse reaction?
        \begin{enumerate}[label=(\alph*)]
          \item Try to predict the answer from what you know.
          \item What does your answer suggest about the
                relative ``reactivity'' of iron and copper?
        \end{enumerate}
\end{enumerate}
\end{problembox}

% ─── Section: Balancing Equations ────────────────────────────────────────────
\subsection{Balancing Equations: The Bookkeeping of Conservation}

If atoms are rearranged but never created or destroyed,
then every chemical equation must \emph{balance}:
the number of each type of atom must be the same on both sides.

This is not just a mathematical exercise.
It is a direct consequence of the law of conservation of mass,
applied at the level of individual atoms.

\subsubsection{The Method}

Consider the combustion of methane (natural gas):

\begin{center}
\ce{CH4 + O2 -> CO2 + H2O} \quad (unbalanced)
\end{center}

Count atoms on each side:

\begin{center}
\begin{tabular}{lcc}
\toprule
Atom & Left side & Right side \\
\midrule
C & 1 & 1 \\
H & 4 & 2 \\
O & 2 & 3 \\
\bottomrule
\end{tabular}
\end{center}

Hydrogen and oxygen are unbalanced.
We adjust coefficients (never subscripts — changing subscripts
changes the \emph{substance} itself):

\begin{center}
\ce{CH4 + 2O2 -> CO2 + 2H2O} \quad (balanced)
\end{center}

\begin{center}
\begin{tabular}{lcc}
\toprule
Atom & Left side & Right side \\
\midrule
C & 1 & 1 \checkmark \\
H & 4 & 4 \checkmark \\
O & 4 & 4 \checkmark \\
\bottomrule
\end{tabular}
\end{center}

\vspace{0.5em}

\begin{problembox}
Balance the following equations.
Show your atom-counting table for each.

\begin{enumerate}[label=\textbf{\arabic*.}]
  \item \ce{H2 + O2 -> H2O}
  \item \ce{Fe + O2 -> Fe2O3}
  \item \ce{Al + HCl -> AlCl3 + H2}
  \item \ce{C3H8 + O2 -> CO2 + H2O} \quad (propane combustion)
  \item \ce{KMnO4 + HCl -> KCl + MnCl2 + H2O + Cl2}
        \quad (this one is harder — try systematically)
\end{enumerate}
\end{problembox}

% ─── Section: Energy in Reactions ────────────────────────────────────────────
\subsection{A First Look at Energy: Why Reactions Happen}

Conservation of mass tells us \emph{what} is rearranged.
It does not tell us \emph{why} rearrangements happen,
or why some reactions release heat and others absorb it.

We can observe:

\begin{itemize}[leftmargin=2em]
  \item Burning wood releases heat (and light):
        the reaction is \newterm{exothermic}.
  \item Dissolving ammonium nitrate in water cools the solution:
        the reaction is \newterm{endothermic}.
  \item Some reactions need a spark or heat to begin
        (like combustion), but once started, sustain themselves.
  \item Other reactions happen spontaneously at room temperature.
\end{itemize}

\begin{center}
\begin{tikzpicture}[
  every node/.style={font=\sffamily\small},
  scale=0.95
]
  % Exothermic energy diagram
  \begin{scope}[xshift=0cm]
    \node[font=\sffamily\small\bfseries, color=JourneyColor]
          at (2, 2.8) {Exothermic Reaction};
    \draw[thick, color=JourneyColor!40, ->] (0,0) -- (4.2,0)
          node[right, font=\sffamily\footnotesize]{reaction progress};
    \draw[thick, color=JourneyColor!40, ->] (0,0) -- (0,2.8)
          node[above, font=\sffamily\footnotesize]{energy};

    % Energy profile curve
    \draw[ultra thick, color=DiscoveryColor, rounded corners=8pt]
          (0.3, 1.8) -- (1.2, 1.8) -- (2.0, 2.5) -- (3.0, 0.7) -- (3.8, 0.7);

    % Labels
    \draw[dashed, color=JourneyColor!50] (0.3, 1.8) -- (-0.05, 1.8)
          node[left, font=\sffamily\footnotesize]{reactants};
    \draw[dashed, color=JourneyColor!50] (3.8, 0.7) -- (4.05, 0.7)
          node[right, font=\sffamily\footnotesize]{products};

    % Arrow showing energy released
    \draw[{Stealth}-{Stealth}, color=WarningColor, thick]
          (3.95, 0.7) -- (3.95, 1.8)
          node[midway, right, font=\sffamily\footnotesize,
               text=WarningColor]{$\Delta H < 0$ heat released};

    % Activation energy
    \draw[{Stealth}-{Stealth}, color=ModuleColor, thick]
          (0.15, 1.8) -- (0.15, 2.5)
          node[midway, left, font=\sffamily\footnotesize,
               text=ModuleColor, align=right]{activation\\energy};
  \end{scope}

  % Endothermic energy diagram
  \begin{scope}[xshift=6.5cm]
    \node[font=\sffamily\small\bfseries, color=JourneyColor]
          at (2, 2.8) {Endothermic Reaction};
    \draw[thick, color=JourneyColor!40, ->] (0,0) -- (4.2,0)
          node[right, font=\sffamily\footnotesize]{reaction progress};
    \draw[thick, color=JourneyColor!40, ->] (0,0) -- (0,2.8)
          node[above, font=\sffamily\footnotesize]{energy};

    % Energy profile curve
    \draw[ultra thick, color=PuzzleColor, rounded corners=8pt]
          (0.3, 0.7) -- (1.2, 0.7) -- (2.0, 2.2) -- (3.0, 1.8) -- (3.8, 1.8);

    % Labels
    \draw[dashed, color=JourneyColor!50] (0.3, 0.7) -- (-0.05, 0.7)
          node[left, font=\sffamily\footnotesize]{reactants};
    \draw[dashed, color=JourneyColor!50] (3.8, 1.8) -- (4.05, 1.8)
          node[right, font=\sffamily\footnotesize]{products};

    % Arrow showing energy absorbed
    \draw[{Stealth}-{Stealth}, color=PuzzleColor, thick]
          (3.95, 0.7) -- (3.95, 1.8)
          node[midway, right, font=\sffamily\footnotesize,
               text=PuzzleColor]{$\Delta H > 0$ heat absorbed};
  \end{scope}
\end{tikzpicture}
\end{center}

\vspace{0.5em}

Why do exothermic reactions exist at all?
Why should rearranging atoms release energy?
This question has a deep answer we will develop fully in Journey~5
when we study thermodynamics.
For now, the intuition is this:
some arrangements of atoms are more \emph{stable} (lower energy)
than others.
When atoms rearrange from a higher-energy to a lower-energy
arrangement, the difference is released as heat.

\begin{puzzlebox}
Consider the reaction:
\begin{center}
\ce{C + O2 -> CO2} \quad ($\Delta H = \SI{-393.5}{\kilo\joule\per\mol}$, exothermic)
\end{center}

\begin{enumerate}[label=(\alph*)]
  \item The reverse reaction (\ce{CO2 -> C + O2}) would absorb the same
        amount of energy. Does this reaction happen spontaneously?
        Why or why not?
  \item Coal is mostly carbon. Burning coal in a power station
        releases energy. Where was that energy stored originally?
        (Think about how coal was formed.)
  \item Diamonds are also pure carbon. They do not burn spontaneously
        at room temperature, even though the reaction is strongly
        exothermic. What does this tell you about the relationship
        between \emph{thermodynamic favourability} (products are
        lower energy) and \emph{whether a reaction actually happens}?
\end{enumerate}
\end{puzzlebox}

% ─── Section: Deep problems ───────────────────────────────────────────────────
\subsection*{Deep Problems}

\begin{deepbox}
\begin{enumerate}[label=\textbf{D\arabic*.}]

  \item \textbf{The Phlogiston Autopsy.}\\
        Phlogiston theory explained combustion, rusting, and
        smelting coherently within one framework —
        and was believed by some of the best scientists in Europe
        for nearly a century.

        \begin{enumerate}[label=(\alph*)]
          \item List three pieces of experimental evidence
                that phlogiston theory could \emph{not} explain
                without adding ad hoc modifications.
          \item Lavoisier did not prove phlogiston didn't exist.
                He showed a different framework explained more,
                more simply.
                What is the philosophical principle here?
                (Look up \emph{Occam's razor} and \emph{parsimony}.)
          \item Is it possible that our current model of atoms
                is also wrong in some important way?
                What would it take to replace it?
        \end{enumerate}

  \item \textbf{Mass Conservation and Nuclear Reactions.}\\
        In a chemical reaction, mass is conserved.
        But in a nuclear reaction — such as the fusion of hydrogen
        in the Sun — mass is \emph{not} conserved: a tiny amount
        of mass is converted to enormous energy, following
        $E = mc^2$.

        \begin{enumerate}[label=(\alph*)]
          \item Why does mass conservation hold for chemical reactions
                but not nuclear ones?
                (Hint: what is the origin of chemical energy
                vs nuclear energy?)
          \item When coal burns and releases heat,
                does the total mass of reactants exactly equal
                the total mass of products, or is there a tiny discrepancy?
                Estimate the discrepancy for \SI{1}{\kilo\gram} of coal
                releasing approximately \SI{3e7}{\joule}.\\
                (Use $E = mc^2$ with $c = 3\times10^8\ \si{\metre\per\second}$.)
          \item Is the discrepancy in (b) measurable with any
                realistic balance?
                What does this tell us about why
                ``mass is conserved in chemical reactions''
                is an excellent approximation even if
                not strictly exact?
        \end{enumerate}

  \item \textbf{Designing the Crucial Experiment.}\\
        Imagine it is 1770.
        You have a precise balance and glassware.
        You suspect phlogiston theory is wrong,
        but you cannot appeal to atomic theory
        (which hasn't been invented yet).

        Design a series of three experiments that would,
        together, make a compelling quantitative case against
        phlogiston theory and in favour of Lavoisier's oxygen theory.

        For each experiment, state:
        \begin{itemize}[nosep]
          \item What you would do
          \item What measurements you would take
          \item What result phlogiston theory predicts
          \item What oxygen theory predicts
          \item Why the result would distinguish the two theories
        \end{itemize}

\end{enumerate}
\end{deepbox}

% ─── Reflection ──────────────────────────────────────────────────────────────
\subsection*{Reflection}

\begin{reflectionbox}
\begin{enumerate}[label=\textbf{R\arabic*.}]
  \item We started this module with a wrong theory (phlogiston)
        that explained much of what was observed.
        \textbf{What made Lavoisier's approach different
        from his predecessors — not in the facts he observed,
        but in the way he approached experiments?}

  \item The law of conservation of mass was derived from
        macroscopic measurements (weighing substances).
        We have not yet mentioned atoms.
        \textbf{What additional hypothesis is needed to go
        from ``mass is conserved'' to ``atoms are rearranged''?}
        Is that hypothesis proven by the weighing experiments?

  \item When we draw an energy diagram for a reaction,
        we show reactants and products at different energy levels.
        \textbf{What is this ``energy'' actually stored in?}
        (Write your current best guess. You will revisit this
        in Journey~4 when we study chemical bonds.)

  \item In 1774, Priestley discovered oxygen but called it
        ``dephlogisticated air.''
        Lavoisier later named it correctly and built
        a new theory around it.
        \textbf{What does this tell us about the relationship
        between experimental discovery and theoretical interpretation?}
        Is discovering a new substance the same as understanding it?
\end{enumerate}
\end{reflectionbox}

\vspace{1.5em}
\begin{center}
\begin{tcolorbox}[
  enhanced, colback=DiscoveryColor!8, colframe=DiscoveryColor,
  width=0.88\textwidth, arc=3pt
]
\centering\sffamily\small
\textbf{What we invented in this module.}\\[4pt]
We invented the habit of \emph{weighing everything}.
Lavoisier's revolution was not primarily intellectual — it was methodological.
He insisted that no explanation was acceptable
unless it accounted for every gram that went in and every gram that came out.
That habit — quantitative accountability — is the foundation
of modern experimental science.
We also named one of the deepest empirical laws in chemistry:
mass is conserved.
Atoms were not yet necessary to say this.
But as we will discover in Module~1.4,
the conservation of mass
\emph{demands} that something like atoms must exist.
\end{tcolorbox}
\end{center}

\vspace{2em}
\noindent\textcolor{JourneyColor!40}{\rule{\textwidth}{0.6pt}}
\begin{center}
\textit{\sffamily\small
Module 1.2 complete. Continue to Module 1.3:
\textbf{Elements and Compounds --- The First Classification}.}
\end{center}

\end{document}